% Pursuing Stacks (cotes 134-2 to 134-8): archive cote 134-3, site page 67 (PDF page 68), typescript page 67.
% typescript. Draft compared: ps2.tex:3472-3530.
% AI-transcribed by Claude Opus 5.5 (claude-opus-5-5, adaptive thinking, effort high), Message Batches API, batch ps1 (msgbatch_01Rv5hCaNTtUt4JHx1yqoHPL), 2026-11-12; not yet checked by eye.
\scanpage{134-3/67}
\mspage{67}

aspheric over $e_{A^\wedge}$ ,  or (what now amounts to the same) that $L_A$ is aspheric.
However,  when faced with the question to decide whether a given A does
indeed \xout{givexrisexto} satisfy  T 2),  it will be convenient to use a \xout{homotopy}
system \xout{\unsure{inxAt}} $(I,e_0,e_1)$ in $A^\wedge$ of which we know beforehand it is a homotopy
interval (the delicate part of this notion being the asphericity of all
products Ixa \unsure{\over{(}{f}}or a in A).  \xout{\unsure{Thexmax}} As part of the ``story retold'',  I recall
now the most natural geometric assumption which will ensure T 2)  i.e.
total \xout{\unsure{asyeli}} asphericity of $A^\wedge$:
\uline{Proposition}   Let A be a small category,  $(I\xout{x},e_0,e_1)$ a homotopy interval
in $A^\wedge$.  Assume that for any a in A,  there exists a homotopy
\[ h_a : Ixa \longrightarrow a \]
between  $id_a$ and a constant map  $c_a : a \longrightarrow a$  (defined by a section of a
over the final element of $A^\wedge$ i.e. by \inserted{a} \xout{\unsure{anxelementxxarx}} map $e \longrightarrow a$).  Then $A^\wedge$is
totally aspheric i.e. every $a \in Ob\ A$ is aspheric over e  (or,  equivalently,
the product in $A^\wedge$\hand{/}of any two elements a, b in A is aspheric).

This is a particular case of the ``homotopy lemma'' (\unsure{\S}.\hand{\unsure{37}}). \xout{x} In fact
we don't even \inserted[typed]{use}\xout{needx}the condition  $e_0 \cap e_1 = \emptyset_{A^\wedge}$  on the``interval'' $(I,e_0,e_1)$,
but we easily see that this condition follows from the existence of the
homotopy and the sections of any a in A over e\hand{,} \deleted{\xout{unless} except in the
\inserted[typed]{trivial} \xout{only} case when \xout{\unsure{$e_0=e_1$ i.e. I=$e_0$xI}xxxxx,xxandxwhenxmoreoverxexxxx} A is equivalent
to the final category i.e. $A^\wedge$  equivalent to (Ens),  and  $I=e_0=e_1=e$ .
However this trivial case isn't interesting for us,  i.e. it is \uline{not} a test
category,  because the Lawvere element (which corresponds to the two-point
set  0,1  in (Ens)) is definitely not aspheric !  (Compare with the more
general corollary 2 on p.53).}

Before returning to the investigation of specific test categories,
I want to come back on some terminology.  The condition on $A^\wedge$\hand{/}that the Law-%
vere element $L_A$ should be aspheric over e has taken lately considerable
geometric significance,  and merits a name.  I will say from now on that
A is a \uline{test} \uline{category} ,  and that $A^\wedge$is an \uline{elementary} \uline{modelizer} ,  if this
condition is satisfied\hand{,}  and if moreover A is aspheric.  This condition
\margin{X} (which amounts to our former  T 1  + T 3 ) is weaker than what \over{W}{w}e had
lately called a test category,  as we had overlooked so far the fact that
T 1 and T 3  alone already imply the basic requirement about  $i_Aj_A(C) \longrightarrow C$
being a weak equivalence for any C , and hence $W_A^{-1}A^\wedge$ being canonically
equivalent to (Hot).  Th\over{\unsure{is}}{us} it seemed for a while that the only handy condi-%
tions we could get for insuring this requirement were T 1) T 2) T 3) ,
all together (in fact , it turned out later that T 2 \xout{xxTx3} already implies
T 1).  When all three conditions T1 to T3 are satisfied,  I'll now say that
